\problemname{Orienteering}
Springoalla has taken up orienteering, but is honestly not
particularly good at it. In fact, despite the guiding arrows
that have been set up, she gets lost almost every time. The forest she runs
in can be seen as a rectangular grid with $N$ rows and $M$ columns,
with arrows of four different types placed: \verb+^+, \texttt{>},
\texttt{v}, and \texttt{<}.

A period character (\texttt{.}) is used to indicate that a cell has no
arrow. Springoalla
enters at the top-left cell, running to the right. When she
reaches an arrow she automatically changes direction and starts running
in the direction the arrow points. However, it sometimes happens that she misses an
arrow, and instead continues straight past it.

Given a position in the forest, how many arrows must Springoalla have
missed at minimum to end up there? Note that she can never have run outside
the forest and whether she misses an arrow is not affected by whether she has been
to that location before (if she e.g. misses the arrow twice it counts
as two misses). There is always at least one way she can have ended up at the given position.

\section*{Input}

The first line contains four numbers $N$, $M$, $R$, and $C$ ($1 \leq R \leq N \leq 800$, $1 \leq C \leq M \leq 800$), the height and width of the forest, and the position (row and column) where Springoalla ends
up (rows are numbered from $1$ to $N$ and columns from $1$ to
$M$). Then follow $N$ lines with $M$ characters each, describing the forest.
Each character will be either \texttt{.}, \texttt{v} (down),
\verb+^+ (up), \texttt{<} (left), or \texttt{>} (right).

\section*{Output}
Output a single number: the minimum number of arrows Springoalla must have missed.

\section*{Scoring}
Your solution will be tested on a set of test groups.
To earn points for a group, you must pass all test cases in that group.

\noindent
\begin{tabular}{| l | l | l |}
  \hline
  Group & Points & Constraints \\ \hline
  $1$   & $50$       & $N \leq 200$, $M \leq 200$ \\ \hline
  $2$   & $50$       & No additional constraints \\ \hline
\end{tabular}

\section*{Explanation of Samples}
Explanation of sample 1: she has missed all downward arrows except one.

Explanation of sample 2: she has missed the middle arrow twice.

Explanation of sample 3: she misses no arrows.
